\chapter{Role of self-attention in a computer vision architecture}
\label{chapter1b}

\startcontents[chapters]
\printmyminitoc

\section{Introduction}

Humans continue to outperform modern AI systems in their ability to flexibly parse and understand complex visual relations. Prior cognitive neuroscience work suggests that attention plays a key role in humans' visual reasoning ability. 

Attention, a fundamental cognitive process, plays a crucial role in shaping our perception, memory, and decision-making. It allows us to selectively focus on relevant information while filtering out distractions, enabling efficient and adaptive behavior in complex environments. By employing selective attention, we give priority to information that is behaviorally relevant while disregarding surrounding stimulation. In neuroscience, the study of attention has provided valuable insights into the workings of the human brain. Neuroscientists have identified distinct neural networks and mechanisms that govern attentional processes, shedding light on how the brain filters and processes sensory inputs, allocates cognitive resources, and guides behavior. Understanding the neural basis of attention has not only deepened our understanding of human cognition but has also provided inspiration for developing attention models in AI.

Attention can be consciously directed towards spatial or non-spatial properties, also known as feature-based attention. By selectively directing our attention, we possess the ability to intentionally focus on particular aspects of our environment. This may include directing our attention to a specific position in space (spatial attention) or highlighting a specific feature, such as a particular color (feature-based attention). When a location or feature is correctly indicated (valid cue; attentional focusing), it results in improved performance for the subsequent stimulus. Conversely, when the cue is incorrect (invalid cue), performance declines as it necessitates reorienting attention to the unexpected target stimulus \citep{alanallport1971parallel,posner1980orienting,posner1987selective}. The impact of cueing, as observed through the difference between valid and invalid trials, can be observed in the activity modulations of neurons in early visual areas \citep{moran1985selective,treue1996attentional,roelfsema1998object}. This effect is further supported by increased activity modulations in early visual areas as revealed by functional magnetic resonance imaging (fMRI) \citep{o1999fmri,ungerleider2000mechanisms}.

Numerous studies have provided evidence that both spatial and feature attention have a modulatory effect on neuronal responses, resulting in an improved signal-to-noise ratio during the encoding of the attended stimulus. Additionally, it has been reported that both types of attention lead to increased neuronal response magnitudes across multiple visual areas \citep{carrasco2011visual}. Moreover, both spatial and feature attention contribute to enhancing the representation of the attended stimulus by reducing neuronal response variability (often measured using the Fano factor) and pairwise noise correlation \citep{cohen2009attention,mitchell2009spatial}. 

While there are notable similarities between feature and spatial attention, several differences have been observed as well. One prominent distinction is that spatial attention is confined to a specific location within the retinotopic map \citep{womelsdorf2006dynamic}, whereas feature attention impacts processing across the entire visual field \citep{saenz2002global}. Furthermore, the temporal dynamics of sensory neuron modulation differ between the two types of attention \citep{hayden2005time}. In order to investigate the underlying neural mechanism that may account for both similarities and differences between spatial and feature attention, \citet{ni2019neuronal} conducted a study involving trained monkeys performing a direction-change detection task. During the task, neuronal activity was recorded from the medial temporal cortex (MT). By manipulating the direction of the Gabors and the attended location, the researchers created three distinct task variants aimed at measuring the neuronal modulation induced by normalization, spatial attention, and feature attention.

In the realm of artificial intelligence, attention mechanisms have become essential components of cutting-edge machine learning algorithms. Inspired by the principles observed in neuroscience, attention models in AI enable machines to focus on salient features or regions of interest within input data, allowing them to allocate computational resources effectively and improve performance on various tasks. By selectively attending to relevant information, AI systems can extract meaningful patterns, make informed decisions, and exhibit more human-like intelligence. The synergy between attention research in neuroscience and AI has led to significant advancements in both fields. Neuroscientists can validate their theories by testing their predictions on AI models, while AI researchers can leverage the findings from neuroscience to design more biologically plausible and efficient attention models. This bidirectional flow of knowledge and insights has the potential to revolutionize our understanding of attention, cognitive processes, and the development of intelligent systems. By bridging the gap between neuroscience and artificial intelligence, we can unlock new perspectives on attention, fostering a deeper understanding of how attention shapes our cognitive abilities and paving the way for more sophisticated and efficient intelligent systems. 



Attention plays a crucial role in addressing the binding problem, which involves integrating various features of a stimulus into a unified object representation. Extensive research in the cognitive and neuroscience fields underscores the significance of attention in this process. The Feature Integration Theory (FIT) proposed by \citet{treisman1980feature} emphasizes how attention facilitates the binding of features, enabling the formation of coherent object representations. \citet{desimone1995neural} work focuses on the neural mechanisms of selective visual attention, elucidating how attention contributes to feature binding and information integration within the visual system. \citet{reynolds2004attentional} explore attention's impact on visual processing, highlighting its role in feature binding and coordinating neural activity across different brain regions. \citet{treisman1998feature} delves into the intricate relationship between feature binding, attention, and object perception. Additionally, \citet{corbetta2002control} investigate the control of goal-directed and stimulus-driven attention, shedding light on their implications for solving the binding problem. Collectively, these studies underscore the indispensable role of attention in integrating features and effectively addressing the binding problem through both cognitive and neural processes.

Furthermore, other works contribute valuable insights into object vision and the temporal dynamics of attention during visual search tasks. \citet{tanaka1996inferotemporal} focuses on the inferotemporal cortex (IT) and its role in object vision, discussing the neural mechanisms involved in recognizing complex object features. \citet{riesenhuber1999cortical} address the ``binding problem'' by examining cortical models and proposing distributed representations and feature-based attention as potential solutions, challenging the notion that cortical models are limited by this problem. \citet{woodman2003serial} investigate the temporal dynamics of attention during visual search tasks and propose a two-stage model involving serial deployment of attention to select target locations, followed by parallel processing within those locations. Together, these works contribute to our understanding of object recognition, the integration of visual features, and the temporal aspects of attention during visual search tasks, complementing the broader literature on attention and the binding problem in cognitive and neuroscience research.

Addressing the binding problem has been a prominent focus in cognitive and neuroscientific research, with computational approaches providing valuable insights. \citet{hommel1998event} introduces the concept of event files, which propose the automatic integration of stimulus-response episodes in memory. This work presents evidence supporting the automatic binding of stimuli and responses, suggesting the creation of temporary associations to optimize processing efficiency. \citet{lisman2013theta} contribute to the field by discussing the theta-gamma neural code, highlighting the significance of synchronized theta and gamma oscillations in neural communication and information processing. Their review emphasizes the role of the theta phase as a temporal framework for precise encoding and integration of information in various cognitive processes. Additionally, \citet{verguts2017binding} presents a computational model known as "binding by random bursts" to elucidate cognitive control mechanisms. This model posits that cognitive control emerges from dynamic interactions between low-level sensory processing and high-level control processes, where random bursts of neural activity act as a binding mechanism, coordinating information flow and facilitating flexible cognitive control. Lastly, \citet{senoussi2022time} investigate time-based binding as a solution and limitation for flexible cognition. They propose that temporal associations between events are crucial for cognitive processing and the binding of information over time. The discussion explores how time-based binding can enhance cognitive flexibility but also introduces constraints in rapidly changing contexts. Together, these studies contribute to our understanding of the binding problem and offer computational models that shed light on cognitive control and the role of time-based binding in flexible cognition.

In the previous chapter, we discussed a benchmark used to evaluate the abilities of machines and compare them with humans. We did this by systematically assessing the ability of modern deep convolutional neural networks (CNNs) to learn to solve the synthetic visual reasoning test (SVRT) challenge, a collection of 23 visual reasoning problems. Our analysis revealed a novel taxonomy of visual reasoning tasks, which can be primarily explained by the type of relations (same-different (SD) versus spatial-relation (SR) judgments) and the number of relations used to compose the underlying rules.

Consistent with the speculated role of attention in solving the binding problem when reasoning about objects \citep{egly1994covert,roelfsema1998object}, prior work by \citet{kim2018not} has shown that combining CNNs with an oracle model of attention and feature binding (i.e., preprocessing images so that they are explicitly and readily organized into discrete object channels) renders SD tasks as easy to learn by CNNs as SR tasks. Here, we build on this work and introduce CNN extensions incorporating spatial or feature-based attention.  In the first set of experiments, we show that these attention networks learn difficult SVRT tasks with fewer training examples than their \textcolor{black}{non}-attentive (CNN) counterparts but that the different forms of attention help on different tasks.


This experiment raises the question: how do attention mechanisms help with learning different visual reasoning problems? There are at least two possible computational benefits: attention could improve model performance by simply increasing its capacity, or attention could help models learn the abstract rules governing object relationships more efficiently. To adjudicate between these two possibilities, we measured the sample efficiency of ResNets pre-trained on SVRT images so that they only had to learn the abstract rules for each SVRT task. We found that attention ResNets and ResNets pre-trained on the SVRT were similarly sample-efficient in learning new SVRT tasks, indicating that attention helps discover abstract rules instead of merely increasing model capacity.


\section{Experiment 1: Self-attention with ResNet50}
\label{sec:exp2}
We sought to identify computational mechanisms that could help ResNets learn the more challenging SVRT tasks revealed by our novel taxonomy. Attention has classically been implicated in visual reasoning in primates and humans \citep{egly1994covert,roelfsema1998object}. Attentional processes can be broadly divided into \textit{\textbf{spatial}} (e.g., attending to all features in a particular image location) vs. \textit{\textbf{feature-based}} (e.g., attending to a particular shape or color at all spatial positions) \citep{desimone1995neural}. The importance of attention for perceiving and reasoning about challenging visual stimuli has also been realized by the computer vision community. There are now a number of attention modules proposed to extend CNN's -- including spatial  (e.g., \citet{Sharma2015-ow,journals/corr/ChenWCGXN15, yang2016, journals/corr/XuS15a,journals/corr/RenZ16}), feature-based (e.g., \citet{stollenga2014deep, chen2017sca, hu2018squeeze}) and hybrid (e.g., \citet{linsley2018global,woo2018cbam}) approaches. Here, we adapt the increasingly popular Transformer architecture \citep{vaswani2017attention} to implement both forms of attention. These networks, which were originally developed for natural language processing, are now pushing the state of the art in computer vision \citep{Zhu2020-nc,carion2020end,Dosovitskiy2020-iq}. \textcolor{black}{Recent work \citep{Ding2020AttentionOL} has also shown the benefits of such architectures and especially attention mechanisms for solving higher-level reasoning problems}.

Transformers are neural network modules usually consisting of at least one ``self-attention'' module followed by a feedforward layer. Here, we introduced different versions of the self-attention module into ResNets to better understand the computational demands of each SVRT task. Transformers' self-attention is applied to and derived from the module's input. By reconfiguring standard Transformer self-attention, we developed versions capable of allocating either spatial or feature-based attention over the input. Specifically, we created these different forms of attention by reshaping the convolutional feature map input to a Transformer. For spatial attention, we reshaped the $ \mathcal{Z} \in \mathcal{R}^{H, W, C}$ (where $H$ is \textit{height}, $W$ is \textit{width} and $C$ is \textit{number of feature channels}) feature maps to $ \mathcal{Z} \in \mathcal{R}^{C, H * W}$ so that the Transformer's self-attention was allocated overall spatial locations. For feature-based attention, we reshaped the convolutional feature maps to $\mathcal{Z} \in \mathcal{R}^{H * W, C}$, enforcing attention to overall features instead of spatial locations. 

\paragraph{Spatial Attention Module (SAM)}

Our first attention module takes a features map $X \in \mathcal{R}^{d_C \times d_H \times d_W}$ as input, where $d_C$, $d_H$, and $d_W$ respectively  refer to the number of channels, height and width of the map, and outputs a features map $Y$ of the same dimensions.

We flatten the spatial dimensions to obtain $X' \in \mathcal{R}^{d_C \times d_N}$, where $d_N = d_H \times d_W$, and we apply the original multi-head self-attention module from \citet{vaswani2017attention} as follows.

We first apply independent linear mappings of the input $X'$ to obtain three feature maps of dimensions $\mathcal{R}^{d \times d_N}$ for each attention head from a total of $n_H$ heads. For the $i^{th}$ head, these maps are known as the query $Q_i$, the key $K_i$ and the value $V_i$, and are obtained such as:
\begin{align*}
Q_i = W^Q_i . X' \\
K_i = W^K_i . X' \\
V_i = W^V_i . X'
\end{align*}
The mappings are parametrized by three matrices $W^Q_i$, $W^K_i$ and $W^V_i$ of dimensions $\mathcal{R}^{d \times d_C}$ for each head. The symbol $.$ denotes a matrix multiplication.

Then, we apply the scaled dot-product attention \citep{vaswani2017attention} to obtain $n_H$ attention heads of dimensions $\mathcal{R}^{d \times d_N}$ such as:
\begin{align}
    H_i = SoftMax(\frac{Q_i . K_i^T}{\sqrt{d}}) V_i
\end{align}

After, we concatenate all attention heads along the first dimension and apply a linear mapping to obtain $Y' \in \mathcal{R}^{d_C \times d_N}$ such as:
\begin{align}
    Z = W^O . Concat(H_1, ..., H_{n_H}) 
\end{align}
The mapping is parametrized by the matrix $W^O \in \mathcal{R}^{d_C \times d}$.

As commonly done, we have a residual connection before applying a layer normalization \citep{ba2016layer} such as:
\begin{align}
    Y' = LayerNorm(Z + X')
\end{align}

Finally, we unflatten $Y'$ to obtain $Y \in \mathcal{R}^{d_C \times d_H \times d_W}$.

We obtain the best results with a representation space of 512 dimensions ($d=512$) and four attention heads ($n_H=4$).


\paragraph{Features-based Attention Module (FBAM)}

Our second attention module is simply obtained by transposing the channel dimension with the spatial dimensions before applying the same transformations.
In other words, we transpose the input $X'$ into $\mathcal{R}^{d_N \times d_C}$ and transpose the output $Y'$ back into $\mathcal{R}^{d_C \times d_N}$. While SAM models attention over the $d_H*d_W$ regions that compose the input features map, FBAM models attention over the $d_C$ features channels.

We obtain the best results with a representation space of 196 dimensions ($d=196$) and one attention head ($n_H=1$).

\begin{figure}[ht]
    \centering
    \includegraphics[width=.9\linewidth]{images/chapter1/architecture-both} 
    \caption{\label{fig:architecture} Location of the Transformer self-attention modules in our ResNet extensions.}
    
\end{figure}

\textcolor{black}{ We added one spatial or feature-based attention after one of the four residual blocks in a ResNet-50. We placed either form of attention module to a ResNet-50 by choosing the location where the addition of attention yielded the best validation accuracy across the SVRT tasks. Through this procedure, we inserted a spatial attention module after the second residual block and a feature-based attention module after the third residual block (Figure~\ref{fig:architecture}).}

To measure the effectiveness of different forms of attention for solving the SVRT, we compared the accuracy of three ResNet-50 models: one capable of spatial attention, one capable of feature-based attention, and one that had no attention mechanisms (``vanilla'') (Figure~\ref{fig:xy_sa}). Spatial attention consistently improved model accuracy on all tasks across all five dataset sizes that models we used for training. The improvement in accuracy is particularly noticeable for the $SD_1$ cluster. Tasks in this sub-cluster are composed of two rules, which ResNets, without attention, struggled to learn. Attention helps ResNets learn these tasks more efficiently. The improvement is also evident for $SD_2$ and $SR_1$. However, the benefit of attention for $SR_2$ is marginal since ResNets without attention already perform well on these tasks. 


We find that feature-based attention leads to the largest improvements for $SD_1$, especially when training on 5k or 10k examples 

(Figure~\ref{fig:fbsa}). On the other hand, spatial attention leads to the largest improvements for $SD_2$ and $SR_1$. This improvement is pronounced when training on 500 or 1000 examples. Taken together, the differential success of spatial versus feature-based attention reveals that their varying attentional demands can explain the task sub-clusters discovered in our data-driven taxonomy.

\begin{figure}[ht!]
\centering
  \begin{subfigure}{.8\textwidth}
    \centering
    \includegraphics[width=.9\linewidth]{images/chapter1/xy_sa_2} 
    \subcaption{Spatial attention}
  \end{subfigure}
  
  \begin{subfigure}{.8\textwidth}
    \centering
    \includegraphics[width=.9\linewidth]{images/chapter1/xy_fb_2}
  \subcaption{Feature-based attention}
 \end{subfigure}%
 \caption{Test accuracies for a baseline ResNet50 vs. the same architecture endowed with the two forms of attention for each of the twenty-three SVRT tasks when varying the number of training examples. A different axis scale is used for $SR_2$ to improve visibility. \textcolor{black}{These curves are constructed by joining task accuracy for five points representing dataset sizes.} }\label{fig:xy_sa}
\end{figure}

\clearpage

\begin{figure}[t]
\centering
 \includegraphics[width=1\linewidth]{images/chapter1/editedfbsa}
 \caption{Test accuracies for 50-layer ResNets with spatial attention (orange), feature-based attention (tan), or no attention (green). Each bar depicts performance after training from scratch on 10k samples. }
 \label{fig:fbsa}
\end{figure}

\begin{figure*}[ht!]
\centering
  \includegraphics[width=1\linewidth]{images/chapter1/slope}
 \caption{The benefit of attention in solving the SVRT is greatest in data-limited training regimes. The x-axis depicts the number of samples for training, and the y-axis depicts a ratio of the average performance of models with attention to models without attention. When the ratio is greater than 1, it shows that attention helps vs. hurts when lower than 1. This gives us five ratios per task and attention process corresponding to each dataset size. We performed a linear fitting procedure for these points and calculated the corresponding slope. This slope characterizes the relative benefits of attention for that particular task as the number of training examples available increases. If the benefit of attention is most evident in lower training regimes, one would expect a relatively small slope. If the benefit of attention is most evident in higher training regimes, one would expect a large slope.
 }\label{fig:slope}
\end{figure*}

To better understand how the ResNet-derived taxonomy found in Experiment 1 can be explained by the need for spatial and feature-based attention, we measured the relative improvement of each form of attention over the vanilla ResNet. For each attention model and task, we calculated the ratio of the test accuracies between the model and the vanilla ResNet50. We repeated this for every training dataset size, then fit a linear model to these ratios to calculate the slope across dataset sizes (see Figure~\ref{fig:slope} for representative examples). We then repeated this procedure for all twenty-three tasks to produce two 23-dimensional vectors containing slopes for each model and every task.

We next used these slopes to understand the attentional demands of each SVRT task. We did this through a two-step procedure. First, we applied a principal component analysis (see Figure~\ref{fig:pca}) to the vanilla ResNet performance feature vectors ($N=15$) derived from Experiment 1. Second, we correlated the principal components with the slope vectors from the two attention models. We restricted our analysis to the first two principal components, which captured  $\sim 93\%$ of the variance in the vanilla ResNet's performance (Figure~\ref{fig:pca}). This analysis revealed a dissociation between the two forms of attention: feature-based attention was most correlated with the first principal component, and spatial attention with the second principal component. Additionally, along the first principal component, we found the broader dichotomy of these 23 tasks into $SD$ and $SR$ clusters, whereas the second principal component divulges the tasks which responded better with spatial attention from tasks requiring either no attention or feature-based attention (as seen in dotted red line along both the axis in Figure~\ref{fig:pca}).  The corresponding Pearson coefficient $r$ and $p$ values are given in Table~\ref{tab:rp}. 

\begin{figure*}[t]
\centering
  \includegraphics[width=.8\linewidth]{images/chapter1/pca_3}
 \caption{Principal component analysis of the twenty-three tasks using the 15-dimensional feature vectors derived from Experiment 1 representing the test accuracy obtained for each task for different dataset sizes and ResNets of varying depths (18, 50 \& 152). The dotted red line represents 4 different bins in which these tasks can be clustered.} \label{fig:pca}
\end{figure*}

\begin{table}[ht]
\centering
\begin{center}
    \caption{Pearson coefficient ($r$) and corresponding $p$ values obtained by correlating the slope vectors of the spatial attention and the feature-based attention modules with the two principal components of Figure~\ref{fig:pca}. See text for details.}
    \label{tab:rp}

\end{center}

\begin{tabular}{lcccc} 
\hline 
   & \multicolumn{2}{c}{$Spatial$} & \multicolumn{2}{c}{$Feature$} \\ \hline 
   & \textbf{r }      & \textbf{p }     & \textbf{r }         & \textbf{p }        \\  \hline \hline 
$PC_{1}$ & 0.466                 & 0.0249                & \textbf{0.649}            & 0.0008                  \\ \hline
$PC_{2}$ & \textbf{-0.652}       & 0.0007               & -0.491                   & 0.0174 \\ \hline
\end{tabular}

\end{table}

To summarize our results from Experiment 2,  we have found that the task clusters derived from ResNet test accuracies computed over a range of depth and training set sizes can be explained in terms of attentional demands. Here, we have shown that endowing these networks with attentional mechanisms helps them learn some of the most challenging problems with far fewer training examples. We also found that the relative improvements obtained over standard ResNets with feature-based and spatial attention are consistent with the taxonomy of visual reasoning tasks found in Experiment 1. More generally, our analysis shows how the relative need for feature vs. spatial attention seems to account for a large fraction of the variance in computational demand required for these SVRT tasks defined in Experiment 1 according to their learnability by ResNets. 

\section{Experiment 2: Feature vs. rule learning}

The learnability of individual SVRT tasks reflects two components: the complexity of the task's visual features and, separately, the complexity of the rule needed to solve the task. To what extent are our estimates of learnability driven by either of these components? We tested this question by training a new set of ResNets without attention according to the procedure laid out in Experiment 1, but with different pre-training strategies. One of the ResNets was pre-trained to learn visual statistics (but not rules) of SVRT images, and another was pre-trained on ImageNet, \citep[a popular computer vision dataset containing natural object categories;][]{deng2009imagenet}.


For pre-training on SVRT, we sampled 5,000 class-balanced images from each of the 23 tasks (5,000 $\times$ 23 = 115,000 samples in total). To ensure the networks did not learn any of the SVRT task rules, we shuffled images and binary class labels across all twenty-three problems while pre-training the network. We then trained models with binary cross-entropy to detect positive examples \textit{without discriminating tasks}. \textcolor{black}{Our assumption is that shuffling images and labels removes any semantic information between individual images and SVRT rules. However, a network with sufficient capacity can still learn the corresponding mapping between arbitrary images and class labels (even though it cannot generalize it to novel samples). To learn this arbitrary mapping, the network has to be able to encode visual features; but by construction, it cannot learn the SVRT task rule.} When training this model and the ImageNet-initialized model to solve individual SVRT tasks, we froze the weights of the convolutional layers and only fine-tuned the classification layers to solve SVRT problems.

Figure~\ref{fig:xy_c_f} shows a comparison between the different architectures in terms of their test accuracies according to the sub-clusters discovered in Experiment 1. These results first confirm that the SVRT pre-training approach works because it consistently outperforms pre-training on ImageNet (Figure~\ref{fig:xy_img}) or training from scratch. Interestingly, for the $SR_{2}$ sub-cluster, we found that the benefits of pre-training on SVRT go down very quickly as the number of training examples grows. We interpret these results as reflecting the fact that generic visual features are sufficient for the task and that the rule can be learned very quickly (somewhere around 500 and 5,000 samples). For $SR_{1}$ sub-cluster, the benefits of starting from features learned from SVRT are somewhat more \textcolor{black}{evident} in low training regimes. Still, these advantages quickly vanish as more training examples are available (the task is learned by all architectures within 5,000 training samples). 

For $SD_{1}$ while there appears to be a \textcolor{black}{noteworthy} advantage of pre-training on SVRT over ImageNet pre-training and training from scratch, the tasks never appear to be fully learned by any of the networks even with 15,000 training examples. This demonstrates the challenge of learning the rules associated with this sub-cluster beyond simply learning good visual representations. Finally, our results also show that the performance gap across all the architectures for $SD_{2}$ vs. $SD_{1}$ increases rapidly with more training examples -- demonstrating the fact that the abstract rule for $SD_{2}$ tasks are more rapidly learned than for $SD_{1}$.

\begin{figure*}[htbp]
\centering
    \includegraphics[width=1\linewidth]{images/chapter1/xy_c_f_2}
 \caption{Test accuracies for a baseline ResNet50 trained from scratch (``No initialization'') vs. the same architecture pre-trained on an auxiliary task in order to learn visual representations that are already adapted to the SVRT stimuli for different \textcolor{black}{numbers} of training examples. The format is the same as used in Figure \ref{fig:xy_sa}. A different axis scale is used for $SR_2$ to improve visibility. \textcolor{black}{These curves are constructed by joining task accuracy for five points representing dataset sizes.}
 }\label{fig:xy_c_f}
\end{figure*}

\textcolor{black}{Finally}, we carried out a similar analysis with the pre-trained network as done in Experiment 2: We built test accuracy vectors for the SVRT pre-trained network trained using all five dataset sizes (.5k, 1k, 5k, 10k, 15k) and searching over a range of optimal learning rates (\textit{1e-4, 1e-5, 1e-6}). This led to a five-dimensional vector, which we  normalized by dividing each entry with the corresponding test accuracy of a baseline ResNet50 trained from scratch. Hence, the normalized vector represents the improvement (ratio larger than 1) or reduction in accuracy (ratio smaller than 1) that results from the pre-training on SVRT for that particular task and training set size. We then calculated the slope vector in $\mathcal R^{(23)}$, which we correlated with the corresponding spatial and feature-based attention vectors from Experiment 2.

\textcolor{black}{We found that task improvements due to SVRT pre-training correlated more strongly with task improvements due to spatial ($r=0.90$, $p=4e-9$) than feature-based attention ($r=0.595$, $p=0.002$)}. This suggests that the observed improvements in accuracy derived from spatial attention are more consistent with learning better feature representations compared to feature-based attention. 

To summarize, in Experiment 3, we have tried to address the question of the learnability of SVRT features vs. rules. We found that  using an auxiliary task to pre-train the networks on the SVRT stimuli in order to learn visual representations beforehand provides learning advantages to the network compared to a network trained from scratch. 

We also found a noteworthy correlation between the test accuracy vector of a network pre-trained on SVRT visual statistics and a similar network endowed with spatial attention. This suggests that spatial attention helps discover the abstract rule more so that it helps improve learning good visual representations for the task. 

\section{Conclusion}

Earlier, \citet{kim2018not} hypothesized that such straining by convolutional networks is due to their lack of attention mechanisms to allow the explicit binding of image regions to mental objects. A similar point was made by \citet{greff2020binding} in the context of the contemporary neural network failure to carve out sensory information into discrete chunks which can then be individually analyzed and compared (see also \citet{10.1007/978-3-540-75555-5_15} for a similar point). Interestingly, this prediction was recently tested using human EEG by \citet{AlamiaENEURO.0267-20.2020} who showed that indeed the brain activity recorded during SD tasks is compatible with greater attention and working memory demands than SR tasks. At the same time, that CNNs can learn SR tasks more efficiently than SD tasks does not necessarily mean that human participants can solve these tasks without attention. Indeed, \citep{logan1994spatial} has shown that SR tasks such as judging insideness require attention under some circumstances.

To assess the role of attention in visual reasoning, we used Transformer modules to endow deep CNNs with spatial and feature-based attention. The relative improvements obtained by the CNNs with the two forms of attention varied across tasks. Many tasks reflected a larger improvement for spatial attention, and a smaller number benefited from feature-based attention. Further, we found that the patterns of relative improvements accounted for much of the variance in the space of SVRT tasks derived in Experiment 1. Overall, we found that the requirement for feature-based and spatial attention accounts well for the taxonomy of visual reasoning tasks identified in Experiment 1.  \textcolor{black}{Our} computational analysis also lead to testable predictions for human experiments by suggesting tasks that either benefit from spatial attention (task \textit{22}) or from feature-based attention (task \textit{21}), tasks that benefit from either form of attention (task \textit{19}), and tasks that do not benefit from attention (task \textit{2}).


Finally, our study has focused on the computational benefits of spatial and feature-based attention for visual reasoning. Future work should consider the role of other forms of attention, including object-based attention \citep{egly1994covert} for visual reasoning.

In our second experiment, we studied the learnability of SVRT features vs. rules. We did this by pre-training the neural networks on auxiliary tasks in order to learn SVRT features before training them to learn the abstract rules associated with individual SVRT problems. Our pre-training methods led to networks that learn to solve the SVRT problems better than networks trained from scratch as well as networks that were pre-trained to perform image categorization on the ImageNet dataset. We have also found that such attention processes seem to contribute more to rule learning than to feature learning. For $SR_1$ sub-cluster we find this type of pre-training to be advantageous in lower training regimes but the benefits rapidly fade away in higher training regimes. In contrast, this pre-training does not allow the tasks from the $SD_1$ sub-cluster to be learned even with 15k samples -- suggesting that the key challenge with these tasks is not to discover good visual representations but rather to discover the rule. This suggests the need for additional mechanisms beyond those implemented in ResNets. This is also consistent with the improvements observed for these tasks with the addition of attention mechanisms. 

In summary, our study compared the computational demands of different visual reasoning tasks. While our focus has been on understanding the computational benefits of attention  and feature learning mechanisms, it is clear that additional mechanisms will be required to fully solve all SVRT tasks. These mechanisms are likely to include working memory which is known to play a role in SD tasks \citep{AlamiaENEURO.0267-20.2020}. Overall, this work illustrates the potential benefits of incorporating brain-like mechanisms in modern neural networks and provides a path forward to achieving human-level visual reasoning. 